\begin{abstract}
開放域表格問答（OTQA）需要從表格集合中檢索相關表格，並透過推理回答自然語言查詢。然而，既有方法主要在小型表格上進行評估，對大型表格的檢索與推理仍缺乏系統性研究。為填補此缺口，本研究基於 DataBench 建立開放域大型表格問答（OLTQA）的評估設定，並提出結合表格結構層級表示與互動式多表推理的方法。在 DataBench 上，本方法在四種大型語言模型下皆取得優於所比較 OTQA 方法的端到端答案正確率，同時維持相對較低的模型輸入量。消融與分析實驗進一步探討候選表格數量、表格表示與互動式表格存取對系統表現的影響。此外，我們在真實金融資料上進行評估，探討此方法在實務應用中的挑戰，並指出未來研究方向。
\end{abstract}

\begin{abstract*}
Open-Domain Table Question Answering (OTQA) requires retrieving relevant tables from a corpus and reasoning over the retrieved tables to answer natural-language queries.
Existing methods are evaluated primarily on small tables, leaving retrieval and reasoning over large tables underexplored.
To address this gap, we derive an evaluation setting from DataBench for Open-Domain Large-Table Question Answering (OLTQA).
We propose a pipeline that combines schema-level table representations with interactive multi-table reasoning.
On DataBench, our method outperforms the evaluated OTQA baselines in end-to-end answer accuracy across four LLM backbones while maintaining competitive input-token usage.
Ablation studies and analyses examine the effects of candidate size, table representation, and interactive table access on system performance.
We further evaluate our method on a real-world financial corpus to investigate challenges arising in practical applications and identify directions for future research.
\end{abstract*}